\chapter{可信评估：从数据划分到指标选择}
\label{ch:evaluation}

评估不是给模型贴上一个分数，而是回答一个具体问题：
模型面对怎样的新数据，完成什么任务，哪些错误最值得关注？
本章先划清数据用途，再用同一组计数解释分类指标，最后处理回归误差与决定系数。

\begin{keyidea}[先定规则，再看结果]
在比较模型前确定数据划分、主要指标和比较流程。
模型、预处理、分类阈值与停止时刻都属于完整方案的一部分。
测试集用于评估已确定方案，不能反复参与这些选择。
\end{keyidea}

\section{训练、验证与测试各自负责什么？}
\label{sec:eval-split}

训练集（Training Set）用于拟合参数；验证集（Validation Set）用于比较候选方案；
测试集（Test Set）在方案确定后用于最终评估。三者按\textbf{数据用途}区分，
不是按程序运行时属于“训练阶段”还是“预测阶段”命名。

\begin{figure}[htbp]
\centering\begin{tikzpicture}[node distance=8mm and 9mm]
\node[box,text width=26mm] (train) {训练集\\拟合参数};
\node[box,text width=30mm,right=of train] (val) {验证集\\比较候选方案};
\node[box,text width=31mm,right=of val] (freeze) {确定完整方案\\必要时重新拟合};
\node[box,text width=25mm,right=of freeze,draw=Amber,fill=Amber!5] (test) {测试集\\最终评估};
\draw[flow] (train)--(val);
\draw[flow] (val)--(freeze);
\draw[flow] (freeze)--(test);
\draw[flow] (val.south) -- ++(0,-9mm) -| (train.south);
\node[lab,below=13mm of val] {验证可以指导方案调整，测试结果不进入这个选择循环};
\end{tikzpicture}

\caption{验证形成选择循环，测试不进入这个循环。图中的箭头表示工作流程，不表示将数据混合。}
\label{fig:eval-split}
\par\smallskip
\readfigure{候选模型在训练集上拟合，然后用验证表现调整方案。
确定方案后，可以按预定流程使用开发数据重新拟合，再对保留的测试集评估。
不能因为测试结果不够好就继续调阈值或换模型，再把同一测试分数作为独立评估。}
\end{figure}

\textbf{训练误差}在训练集上算，\textbf{验证误差}在验证集上算，\textbf{测试误差}在测试集上算。
分类错误率在相应数据集 $D$ 上定义为
\[
 \widehat e_D=\frac1{|D|}\sum_{i\in D}\mathbf1[\hat y_i\ne y_i].
\]
同一公式用于不同数据集时，角色不同。训练错误率低，不等于泛化误差低。

\begin{pitfall}[独立划分不只是“行号不重叠”]
同一设备的多次记录若跨训练与测试，仍可能泄漏设备特征。
面向未来预测时，应按时间顺序划分，不能随意打乱未来信息。
划分应模拟实际使用条件；随机独立同分布划分只是常见设置之一。
\end{pitfall}

\section{交叉验证：评估整个拟合流程}
\label{sec:eval-cv}

$K$折交叉验证（$K$-Fold Cross-Validation）在开发数据内轮流留出一折作验证，
其余 $K-1$ 折训练。每折都从头拟合预处理和模型，再汇总各折验证结果。
\textbf{不从这 $K$ 个模型中简单挑出分数最高的一个作为最终模型。}

\begin{figure}[htbp]
\centering\begin{tikzpicture}[x=22mm,y=9mm]
\foreach \j in {1,...,5}{\node[lab] at (\j,0) {子集\j};}
\foreach \i in {1,...,5}{
 \node[lab,anchor=east] at (.35,-\i) {第\i 折};
 \foreach \j in {1,...,5}{
  \ifnum\i=\j
   \node[draw=Amber,fill=Amber!12,minimum width=20mm,minimum height=7mm,font=\small] at (\j,-\i) {验证};
  \else
   \node[draw=Teal,fill=Pale,minimum width=20mm,minimum height=7mm,font=\small] at (\j,-\i) {训练};
  \fi
 }
}
\node[lab] at (3,-6.2) {每折重新拟合预处理与模型；独立测试集不出现在这张表中};
\end{tikzpicture}

\caption{五折示例：每个样本恰作一次折内验证，其余四次用于折内训练。}
\label{fig:eval-cv}
\par\smallskip
\readfigure{每一行是一次独立拟合，橙色格仅用于该折验证。
把每个候选方案都按相同划分评估，再依据汇总指标选择方案。
最后按选定方案用开发数据重新拟合，独立测试集继续留在流程之外。}
\end{figure}

\subsection{预处理也必须在折内拟合}
标准化的均值与方差、填补缺失值的统计量、特征筛选规则等，
都只能从该折的子训练集学习。先在全部数据上做这些步骤，再交叉验证，会让验证信息提前进入训练。
把预处理与估计器放在同一拟合流水线中，有助于维持这个边界。

\begin{center}
\small
\begin{tabularx}{\textwidth}{P{33mm}Y}
\toprule
数据情况 & 常见划分考虑\\
\midrule
样本近似独立同分布 & 普通K折或随机留出。\\
分类且类别不均衡 & 可分层，让各折类别比例较稳定；少数类数量仍限制折数。\\
同一对象多次观测 & 按对象分组，避免同一对象跨训练与验证。\\
预测未来 & 按时间向前验证，必要时留间隔防止窗口重叠。\\
\bottomrule
\end{tabularx}
\end{center}

选择折数要权衡样本量和计算开销，不存在“一律10折”的规则。
重复用同一套交叉验证调参仍可能产生选择偏差；可用独立测试集，或在适当设置下用嵌套交叉验证评估选择流程。

\section{混淆矩阵：先声明谁是正类}
\label{sec:eval-confusion}

二分类的正类（Positive Class）是需要明确指定的类别，不代表“好”的结果。
例如设备故障检测中，可以把“故障”设为正类。
混淆矩阵（Confusion Matrix）把真实标签与预测标签交叉计数。

\begin{figure}[htbp]
\centering\begin{tikzpicture}[x=38mm,y=18mm]
\node[lab] at (1.5,2.3) {预测类别};
\node[lab] at (1,1.65) {正类};\node[lab] at (2,1.65) {负类};
\node[lab,rotate=90] at (-.05,.5) {真实类别};
\node[lab] at (.45,1) {正类};\node[lab] at (.45,0) {负类};
\node[box,minimum width=35mm,minimum height=15mm] at (1,1) {TP：18};
\node[box,minimum width=35mm,minimum height=15mm,draw=Amber,fill=Amber!8] at (2,1) {FN：12};
\node[box,minimum width=35mm,minimum height=15mm,draw=Amber,fill=Amber!8] at (1,0) {FP：7};
\node[box,minimum width=35mm,minimum height=15mm] at (2,0) {TN：63};
\end{tikzpicture}

\caption{100个样本的示例。本书约定行是真实类别、列是预测类别，正类都排在前面。}
\label{fig:eval-confusion}
\par\smallskip
\readfigure{30个真实正类中，18个被检出，12个漏检；
70个真实负类中，7个被误报为正类，63个正确判为负类。
因此预测正类有25个，预测负类有75个。}
\end{figure}

\begin{center}
\small
\begin{tabularx}{\textwidth}{P{24mm}Y Y}
\toprule
计数 & 英文与中文 & 本例含义\\
\midrule
TP & True Positive，真正例 & 真实为正，预测也为正：18。\\
FN & False Negative，假负例 & 真实为正，预测为负：12。\\
FP & False Positive，假正例 & 真实为负，预测为正：7。\\
TN & True Negative，真负例 & 真实为负，预测也为负：63。\\
\bottomrule
\end{tabularx}
\end{center}

多分类时可扩展为 $K\times K$ 计数矩阵，主对角线是分类正确的数量。
非对角线告诉我们哪些类别互相混淆，比单个准确率提供更多信息。

\begin{pitfall}
不同软件可能按负类在前或其他标签顺序输出矩阵。读取位置之前先核对轴含义和标签顺序。
如果只写“左上角是TP”而没有声明约定，就容易把四种计数认反。
\end{pitfall}

\section{准确率、查准率与召回率：分母决定问题}
\label{sec:eval-precision-recall}

准确率（Accuracy）问“全部样本中有多少判对”；
查准率（Precision，也称精确率）问“预测为正的样本中有多少真的为正”；
召回率（Recall）问“真实正类中有多少被检出”。
三者的分母分别是全部样本、预测正类、真实正类。

\begin{align}
 \mathrm{Accuracy}&=\frac{TP+TN}{TP+TN+FP+FN},\\
 \mathrm{Precision}&=\frac{TP}{TP+FP},\qquad
 \mathrm{Recall}=\frac{TP}{TP+FN}.
\end{align}

对图~\ref{fig:eval-confusion}，准确率为 $81/100=0.81$，
查准率为 $18/25=0.72$，召回率为 $18/30=0.60$。
分类错误率为 $1-\mathrm{Accuracy}=0.19$。

\subsection{把其它常见名称放回同一组计数}
在二分类正类口径下，灵敏度（Sensitivity）和真正例率
（True Positive Rate, TPR）就是召回率。
特异度（Specificity，也称真负例率TNR）看真实负类中有多少判对；
假正例率（False Positive Rate, FPR）看真实负类中有多少被误报：
\[
 \mathrm{TPR}=\frac{TP}{TP+FN},\qquad
 \mathrm{TNR}=\frac{TN}{TN+FP},\qquad
 \mathrm{FPR}=\frac{FP}{FP+TN}=1-\mathrm{TNR}.
\]
本例特异度为0.90，假正例率为0.10。
\textbf{FPR的分母是真实负类，不是预测正类，因此它不等于 $1-\mathrm{Precision}$。}

\begin{keyidea}[为什么准确率可能不够？]
本例一律预测负类也有70\%准确率，却检不出任何正类。
类别极不均衡时，这种现象更明显。应报告混淆矩阵，并结合漏检和误报的代价选择指标。
没有单个指标能在所有任务上替代这项判断。
\end{keyidea}

若某指标公式的分母为0，其数学定义无法直接使用。
软件可能返回0、缺失值或警告，应记录约定，而不是把这些输出混同于有足够数据支持的测量值。

\section{F1、平均方式与分类阈值}
\label{sec:eval-f1}

F1分数（F1 Score）综合查准率与召回率。
在两者定义有效且其和大于0时，它是调和平均：
\[
 F_1=\frac{2PR}{P+R}
     =\frac{2TP}{2TP+FP+FN},\qquad
 P=\mathrm{Precision},\ R=\mathrm{Recall}.
\]
本例为 $36/(36+7+12)=36/55\approx0.655$。
这里的 $P$ 不是Accuracy；本书避免使用含义不清的“精度”替代指标全名。
F1不直接包含TN，因此也不能代替对真负类表现和错误代价的检查。

\subsection{多分类汇总必须写明平均方式}
可先把每个类别轮流视为正类，求各类Precision、Recall、F1，再决定怎样汇总。

\begin{center}
\small
\renewcommand{\arraystretch}{1.4}
\begin{tabularx}{\textwidth}{P{28mm}Y}
\toprule
方式 & 汇总口径\\
\midrule
宏平均（Macro） & 各类别的指标等权平均，少数类同样占一份权重。\\
加权平均（Weighted） & 按真实类别样本数加权，各类权重可能差别很大。\\
微平均（Micro） & 先累加各类TP、FP、FN，再用总计数计算指标。\\
\bottomrule
\end{tabularx}
\end{center}

在覆盖全部类别的\textbf{单标签多分类}中，micro-Precision、micro-Recall和micro-F1都等于Accuracy。
宏平均F1应当先求各类F1再平均，通常不等于把宏平均Precision和Recall代入F1公式。

\subsection{得分与阈值也是两个对象}
若模型输出正类得分 $s(\vct x)$，可按 $s\geq\tau$ 判为正类。
降低阈值通常会让更多样本判为正类，召回率不会下降，但Precision未必单调变化。
阈值应在验证流程内选择，然后固定用于测试；不能在测试集上挑最高F1再报告为独立测试结果。

\begin{pitfall}
F1对Precision和Recall给予特定形式的折中，不是通用的业务收益函数。
当两种错误的代价差异明显时，应明确代价、约束或使用合适的其它指标，而非仅追求F1最大。
\end{pitfall}

\clearpage
\section{多标签：逐个标签正确，还是整组正确？}
\label{sec:eval-multilabel}
一条记录可以同时有多个标签，例如一张图片既有“猫”也有“室内”。
用$y_{ik},\hat y_{ik}\in\{0,1\}$表示第$i$条记录、第$k$个标签的真实与预测状态，
共$n$条记录、$K$个标签。这与每条记录只选一个类别的多分类不同。

\begin{figure}[H]
\centering\begin{tikzpicture}[x=12mm,y=8mm]
\node[lab] at (1.5,4) {真实标签};\node[lab] at (7.5,4) {预测标签};
\foreach \x/\name in {0/A,1/B,2/C,3/D}{
 \node[lab] at (\x,3) {\name};\node[lab] at (\x+6,3) {\name};}
\foreach \y in {0,1,2}{\node[lab,anchor=east] at (-.6,\y) {样本\pgfmathparse{int(3-\y)}\pgfmathresult};}
\foreach \x/\y/\v in {0/2/1,1/2/0,2/2/1,3/2/0,0/1/0,1/1/1,2/1/0,3/1/0,0/0/1,1/0/1,2/0/0,3/0/0}{
 \node[draw=Teal,fill=Pale,minimum width=11mm,minimum height=7mm,font=\small] at (\x,\y) {\v};}
\foreach \x/\y/\v in {6/2/1,7/2/0,9/2/0,6/1/0,7/1/1,8/1/0,6/0/1,7/0/1,8/0/0,9/0/0}{
 \node[draw=Teal,fill=Pale,minimum width=11mm,minimum height=7mm,font=\small] at (\x,\y) {\v};}
\foreach \x/\y/\v in {8/2/0,9/1/1}{
 \node[draw=Amber,fill=Amber!15,minimum width=11mm,minimum height=7mm,font=\small] at (\x,\y) {\v};}
\draw[flow] (3.8,1)--(5.1,1);
\end{tikzpicture}

\caption{三个样本、四个标签；橙色格为两处错误。每一行是一个标签集合。}
\end{figure}
\readfigure{第1条漏掉标签C，第2条多报标签D，第3条完全正确。
12个二元位置错2个，但只有1条记录的整组标签完全匹配。}

\begin{align}
 L_{\mathrm{Hamming}}&=\frac1{nK}\sum_{i=1}^n\sum_{k=1}^K
      \mathbf1[y_{ik}\ne\hat y_{ik}],\\
 \mathrm{SubsetAccuracy}&=\frac1n\sum_{i=1}^n
      \mathbf1[\vct y_i=\hat{\vct y}_i].
\end{align}
本例Hamming损失为$2/12=1/6$，越小越好；整组匹配准确率为$1/3$，越大越好。
整组0--1损失为$1-\mathrm{SubsetAccuracy}=2/3$，不是Hamming损失。
这些定义不要求不同标签在统计上独立。

\subsection{F1也要说明在哪个方向汇总}
各标签先用各自计数求F1，再等权平均，是Macro-F1；先将所有标签的计数相加再求F1，是Micro-F1。
也可先对每条记录的标签集合求F1，再对记录平均，称Samples-F1。
本例总TP=4、FP=1、FN=1，Micro-F1为$8/10=0.8$；
各标签F1为$1,1,0,0$，Macro-F1为0.5；Samples-F1为$(2/3+2/3+1)/3=7/9$。

\begin{pitfall}[很多负标签会让Hamming看起来很小]
若100个标签中每条记录恰有1个正标签，全部预测0的Hamming损失只有0.01，
整组匹配准确率和召回率却为0。应结合正标签表现、各标签支持数与实际错误代价。
空集合导致某些F1分母为0时，另行说明软件约定。
\end{pitfall}

输出通常使用每标签Sigmoid和二元损失；各标签阈值在验证流程中选择。
损失按标签求和还是平均也需写明，不能直接比较不同缩放下的数值。


\section{ROC：让阈值变化起来}
\label{sec:eval-roc}

ROC曲线（Receiver Operating Characteristic Curve）汇总不同阈值下的分类表现。
本书采用常见坐标：\textbf{横轴FPR，纵轴TPR}，两轴都从0到1。
它需要模型输出连续或有序得分，不要求这些得分已是校准概率。

\begin{figure}[htbp]
\centering\begin{tikzpicture}
\begin{axis}[width=.78\textwidth,height=69mm,xmin=0,xmax=1,ymin=0,ymax=1,
 xlabel={假正例率 FPR},ylabel={真正例率 TPR},xtick={0,.25,.5,.75,1},ytick={0,.25,.5,.75,1},
 tick label style={font=\small},grid=major,grid style={gray!15}]
\addplot[draw=none,fill=Teal!10] coordinates {(0,0)(0,.25)(.25,.5)(.5,.5)(.5,.75)(.75,1)(1,1)(1,0)} \closedcycle;
\addplot[Teal,thick,mark=*] table[x=fpr,y=tpr,col sep=comma]{figures/evaluation/roc.csv};
\addplot[Quiet,dashed,domain=0:1]{x};
\node[font=\small,text=Ink] at (axis cs:.68,.23) {AUC $=0.6875$};
\end{axis}
\end{tikzpicture}

\caption{八个样本的ROC示例，得分在下一节给出。并列得分作为一组改变阈值。虚线为对角参考线。}
\label{fig:eval-roc}
\par\smallskip
\readfigure{阈值高于全部得分时，没有正类预测，对应 $(0,0)$。
逐步降低阈值，会逐组纳入得分相同的样本；阈值低到所有样本都判为正时到达 $(1,1)$。
曲线越靠近左上角，表示能在较低假正例率下取得较高真正例率。}
\end{figure}

每一个曲线点都代表一套阈值选择，不是同一个分类器在固定阈值下同时达到的多种结果。
若实际只能接受很低的误报率，应重点看相应区域，而不是只比较整个曲线面积。

\begin{pitfall}[不要把Specificity直接当作本图横坐标]
有些领域会使用从1向0排列的特异度横轴，等价展示同一关系。
必须同时声明坐标含义和方向，不能把“从0到1的特异度轴”误画成标准FPR轴。
计算ROC与AUC时，评估数据还需包含正、负两类。
\end{pitfall}

\section{AUC：排序能力与并列得分}
\label{sec:eval-auc}

ROC曲线下面积称为ROC-AUC。对经验二分类得分，一种直观算法是逐对比较正、负样本：
\begin{equation}
 \widehat{\mathrm{AUC}}=
 \frac{\displaystyle\sum_{i\in+}\sum_{j\in-}
 [\mathbf1(s_i>s_j)+\tfrac12\mathbf1(s_i=s_j)]}{n_+n_-}.
 \label{eq:eval-auc}
\end{equation}
其中 $n_+,n_-$ 是正、负样本数，$s$ 越大表示越倾向正类。
正类得分更高记1分，相同记半分，更低记0分。

\begin{center}
\small
\begin{tabularx}{.9\textwidth}{P{32mm}Y}
\toprule
真实类别 & 模型得分\\
\midrule
正类，4个 & $0.9,\ 0.8,\ 0.5,\ 0.4$\\
负类，4个 & $0.8,\ 0.6,\ 0.4,\ 0.1$\\
\bottomrule
\end{tabularx}
\end{center}

四个正类与所有负类逐一比较，分别得到 $4,3.5,2,1.5$ 分，
所以 $\widehat{\mathrm{AUC}}=11/16=0.6875$。
这也等于图~\ref{fig:eval-roc} 按并列分组后作线性连接的梯形面积。
若在并列组内部任意排列正负样本再连接，可能得到受排列影响的错误面积。

\subsection{AUC告诉了什么，又没告诉什么？}
AUC衡量得分把正类排在负类之前的能力，不直接衡量概率校准，也不指定实际部署阈值。
严格递增地变换全部得分不会改变其排序，因而不会改变相应AUC。
一个固定阈值下的准确率却可能改变。

\begin{keyidea}
AUC为0.5只描述这一排序汇总量，不能断言模型的全部行为与随机猜测相同。
AUC为1表示在当前评估样本上正类得分都高于负类，也不是对未来数据的完美保证。
报告AUC时，应同时说明数据划分与适用任务；多分类还要说明如何组合类别。
\end{keyidea}

\section{类别比例变化时，Precision也会变}
\label{sec:eval-prevalence}

ROC使用类内比例TPR和FPR；Precision还取决于总体中正类有多常见。
记正类比例为 $\pi$，在类条件表现保持不变的假设下，
\begin{equation}
 \mathrm{Precision}=
 \frac{\pi\,\mathrm{TPR}}
 {\pi\,\mathrm{TPR}+(1-\pi)\,\mathrm{FPR}}.
 \label{eq:eval-prior}
\end{equation}
例如保持 $\mathrm{TPR}=0.8$、$\mathrm{FPR}=0.1$：

\begin{center}
\small
\begin{tabularx}{.9\textwidth}{P{32mm}Y Y}
\toprule
正类比例 & 真正例占全部样本比例 & Precision\\
\midrule
50\% & $0.5\times0.8=0.4$ & $0.4/(0.4+0.05)\approx0.889$\\
1\% & $0.01\times0.8=0.008$ & $0.008/(0.008+0.099)\approx0.075$\\
\bottomrule
\end{tabularx}
\end{center}

两种情形的TPR与FPR相同，预测正类的可靠程度却相差很大。
因此，把在均衡抽样数据上得到的Precision直接当作实际低发生率场景的Precision，通常不可靠。
若类条件分布也变化，上式中的TPR与FPR本身还可能改变。

\subsection{不均衡任务的报告方式}
除Accuracy外，报告正类支持数、Precision、Recall、阈值和混淆矩阵。
Precision--Recall曲线有助于观察查准与查全的关系；若报告其汇总面积或平均精确率，
还要明确计算约定，不能把所有不同定义的“PR面积”当作同一个数。

\begin{pitfall}
“ROC-AUC对类别比例不敏感”依赖类条件得分分布保持不变，
不表示不均衡任务仅看AUC就够了，更不表示样本少的正类能得到精确稳定的估计。
验证/测试集应代表目标使用情形，并同时考虑统计不确定性。
\end{pitfall}

\section{回归误差：先看单位与残差}
\label{sec:eval-regression}

回归任务预测数值。对 $n$ 个样本，记真实目标 $y_i$、预测 $\hat y_i$，
残差 $e_i=y_i-\hat y_i$。常见指标为
\begin{align}
 \mathrm{MAE}&=\frac1n\sum_i|e_i|,\\
 \mathrm{MSE}&=\frac1n\sum_i e_i^2,\qquad
 \mathrm{RMSE}=\sqrt{\mathrm{MSE}}.
\end{align}
MAE与RMSE和目标变量同单位；MSE的单位是目标单位的平方。
平方误差对大残差惩罚更强，但不能因此断言MAE完全不受异常值影响。

\begin{figure}[htbp]
\centering\begin{tikzpicture}
\begin{axis}[width=.83\textwidth,height=60mm,xmin=.7,xmax=3.3,ymin=.5,ymax=4.5,
 xlabel={样本编号 $i$},ylabel={目标与预测值},xtick={1,2,3},ytick={1,2,3,4},
 grid=major,grid style={gray!15},tick label style={font=\small},
 legend style={font=\small,at={(.5,-.23)},anchor=north,legend columns=3,draw=none}]
\addplot[Teal,only marks,mark=*,mark size=2pt] coordinates {(1,1)(2,2)(3,3)};
\addlegendentry{真实 $y_i$}
\addplot[Amber,only marks,mark=square*,mark size=2pt] coordinates {(1,1)(2,2)(3,4)};
\addlegendentry{预测 $\hat y_i$}
\addplot[Quiet,dashed,domain=.7:3.3]{2};
\addlegendentry{目标均值 $\bar y=2$}
\draw[<->,Amber,thick] (axis cs:3,3.08)--(axis cs:3,3.92);
\end{axis}
\end{tikzpicture}

\caption{真实目标为 $(1,2,3)$，预测为 $(1,2,4)$。箭头标出第3个样本的残差大小1；前两个重合点残差为0。}
\label{fig:eval-regression}
\par\smallskip
\readfigure{残差向量为 $(0,0,-1)$，因此MAE与MSE都为 $1/3$，
RMSE为 $\sqrt{1/3}\approx0.577$。
图中水平线是当前评估目标的均值，不是把模型预测都替换成均值。}
\end{figure}

\begin{keyidea}[指标与训练目标可以使用相同公式，但用途不同]
MSE既可以作为训练目标，也可以作为验证或测试指标。
它在测试集上计算时并不会反过来参与训练。
比较模型要使用相同数据划分和相同目标尺度；如果训练时变换了目标，需说明评估在哪个尺度上进行。
\end{keyidea}

\section{R²与平方和：把适用条件写出来}
\label{sec:eval-r2}

决定系数（Coefficient of Determination, $R^2$）把残差平方和与
“始终预测本评估集目标均值”的基准比较。设 $\bar y=n^{-1}\sum_i y_i$：
\begin{align}
 \mathrm{TSS}&=\sum_i(y_i-\bar y)^2, &
 \mathrm{RSS}&=\sum_i(y_i-\hat y_i)^2,\\
 R^2&=1-\frac{\mathrm{RSS}}{\mathrm{TSS}},\qquad & &\mathrm{TSS}>0.
 \label{eq:eval-r2}
\end{align}
完美预测的 $R^2=1$；与该均值基准同样的平方误差给出0；比它更差则为负。
若评估集目标全相同，TSS为0，数学定义的分母失效；软件可能采用约定值，应说明处理方式。
当前目标均值只是这个评估统计量的参照，不是允许用测试标签拟合部署模型；实际可部署的均值预测器应从训练数据估计。

\begin{center}
\small
\begin{tabularx}{\textwidth}{P{29mm}Y Y Y}
\toprule
目标均为 $(1,2,3)$ & 预测 & RSS & $R^2$\\
\midrule
图中模型 & $(1,2,4)$ & $1$ & $0.5$\\
均值基准 & $(2,2,2)$ & $2$ & $0$\\
较差预测 & $(4,4,4)$ & $14$ & $-6$\\
\bottomrule
\end{tabularx}
\end{center}

\subsection{解释平方和不总等于总平方和减残差平方和}
常见定义 $\mathrm{ESS}=\sum_i(\hat y_i-\bar y)^2$。
一般展开为
\[
 \mathrm{TSS}=\mathrm{RSS}+\mathrm{ESS}
   +2\sum_i(y_i-\hat y_i)(\hat y_i-\bar y).
\]
只有交叉项为0时才可直接相加；带截距普通最小二乘的训练集拟合是常见情形。
本例预测 $(1,2,4)$ 的ESS为5，交叉项为 $-4$，所以 $2=1+5-4$，
不能把RSS和ESS无条件相加。

\begin{pitfall}
$R^2$并非对任意预测都等于相关系数平方。
带截距的一元最小二乘训练集拟合、相应方差非0时，有 $R^2=r(x,y)^2$；
任意模型或测试集不能照搬。本节图中预测与目标的相关系数平方约为0.964，实际 $R^2$只有0.5。
\end{pitfall}

\section{速查与自测}
\label{sec:eval-summary}

\begingroup
\small
\renewcommand{\arraystretch}{1.2}
\begin{xltabular}{\textwidth}{P{38mm}Y}
\toprule
要回答的问题 & 先看什么\\
\midrule
\endhead
\bottomrule
\endfoot
模型参数在哪里拟合？ & 训练集；交叉验证中是每折子训练集。\\
阈值或超参数在哪里选？ & 验证流程内；最终测试集保持独立。\\
预测正类有多少是真的？ & Precision：$TP/(TP+FP)$。\\
真实正类检出了多少？ & Recall/TPR：$TP/(TP+FN)$。\\
真实负类误报了多少？ & FPR：$FP/(FP+TN)$。\\
跨阈值的排序表现怎样？ & ROC与ROC-AUC；并列正负得分按半对计分。\\
数值误差有多大？ & MAE、RMSE与残差分布，注意单位。\\
比目标均值基准好多少？ & $R^2=1-\mathrm{RSS}/\mathrm{TSS}$，可能为负。\\

\end{xltabular}
\endgroup

\subsection{三个自测问题}
\begin{enumerate}
\setlength{\itemsep}{1pt}\setlength{\parsep}{0pt}
 \item 若TP=9、FP=1、FN=6、TN=84，Precision与Recall分别是多少？
 \item 两个正类得分为0.8、0.4，两个负类为0.8、0.2，AUC是多少？
 \item 若TSS=10、RSS=15，$R^2$是多少？能否直接认为计算有误？
\end{enumerate}
\textbf{参考答案：}（1）0.9与0.6。（2）$(1.5+1)/4=0.625$。
（3）$-0.5$；它表示平方误差比该均值基准更大，不必然是计算错误。

\par\begingroup\small
\noindent\textbf{进一步查阅：}
scikit-learn：\href{https://scikit-learn.org/stable/common_pitfalls.html}{数据泄漏与预处理常见问题}。
指标文档：\href{https://scikit-learn.org/stable/modules/generated/sklearn.metrics.roc_auc_score.html}{roc\_auc\_score}、
 \href{https://scikit-learn.org/stable/modules/generated/sklearn.metrics.r2_score.html}{r2\_score}，核对输入口径与边界约定。
\par\endgroup
